\subsection{An individually specified finite counterexample}

\begin{claim}{AG.38}{Refuted, with an explicit finite counterexample: nontrivial primitive-cover existence}
\label{clm:AGexplicitprimitivecounterexample}
Let
\[
 \begin{split}
 M(Z)&=Z^{38}+Z^{14}+Z^{10}+Z^9+Z^8+Z^5+Z^2+Z+1,\\
 k&=\mathbb F_2[Z]/(M(Z)),\qquad t=Z\bmod M(Z),
 \qquad d=t^3+1.
 \end{split}                                                    \tag{AG.38.1}
\]
The polynomial \(M\) is irreducible. Let \(C/k\) be the smooth
projective model of the affine equation
\[
             y^2+y=\left(\frac{u^2+t^2u+t}{d}\right)^3,
                                                               \tag{AG.38.2}
\]
and let \(P\) be its unique point at infinity. This is a
supersingular genus-two curve, with \(P\in C(k)\).
For every connected cyclic etale cover
\(D\to C_{\bar k}\) of odd degree \(n>1\), and every generator
\(\sigma\) of its deck group, the abelian subvariety
\[
 P_n(D)=\operatorname{im}\left(
       \prod_{\ell\mid n}(1-\sigma^{n/\ell}):J_D\longrightarrow J_D
                              \right)
\]
has positive two-rank. Thus the nontrivial-cover version of A6 is
false for this completely specified curve, independently of the
counting existence argument.

The degree-one qualification is essential: \(J_C\) has two-rank
zero, so the identity cover satisfies the literal variant that
permits \(n=1\). The example also satisfies A1 and is not a
collision or a counterexample to A1.
\end{claim}

\begin{proof}
First, the standard-library bit-polynomial program
\path{primitive_parameter_field_check.py} verifies
\[
 Z^{2^{38}}\equiv Z\pmod M,
 \qquad
 \gcd(M,Z^{2^2}-Z)=\gcd(M,Z^{2^{19}}-Z)=1.
                                                               \tag{AG.38.3}
\]
Since the prime divisors of \(38\) are \(2\) and \(19\),
the finite-field irreducibility criterion proves that \(M\) is
irreducible. Its retained output gives every successive squaring
residue, using bitwise addition, multiplication, and division in
\(\mathbb F_2[Z]\). In particular \(t\notin\mathbb F_4\),
and \(d\ne0\).

Claim~AG.33 applies to (AG.38.2). It supplies the genus,
supersingularity, unique point at infinity \(P\), which is rational,
scalar Frobenius
\([-2^{19}]\), and the pointed degree-two isogeny
\(\alpha:J_C\to E^2\), where \(E:y^2+y=x^3\). On the
embedded theta curve \(C-P\), its two coordinates have
\[
 A(u)=\frac{u^2+t^2u+t}{d},\qquad
 B(u)=\frac{tu^2+u+t^2}{d}                              \tag{AG.38.4}
\]
as their elliptic \(x\)-coordinates. Both elliptic images of a
finite point are finite, and \(P\) maps to \((0_E,0_E)\).
The theta divisor relevant to the primitive two-rank criterion
is precisely \(C-P\).

By Claims~AG.32, AG.37, and~\ref{clm:ACpacketconductors}, the
only remaining possible conductors of a primitive-zero cover are
\(191\) and \(229\). We prove the stronger assertion that
\(C-P\) has no nonzero point killed by either integer.

Here is a direct characteristic-two verification of the elliptic
torsion polynomials used in that test. On \(E\), the tangent
and chord formulas give the rational-function identities
\[
 x([2]R)=x(R)^4,
 \qquad
 x(R+S)+x(R-S)=\frac{1}{(x(R)+x(S))^2}.
                                                               \tag{AG.38.5}
\]
For the second identity, replacing \(S=(x_S,y_S)\) by
\(-S=(x_S,y_S+1)\) changes the chord slope by
\(1/(x_R+x_S)\); the other terms in the two addition formulas
cancel. The tangent slope is \(x(R)^2\), giving the first
identity. Hence the rational functions \(r_n(X)=x([n]R)\),
where \(X=x(R)\), satisfy
\[
 r_1=X,\qquad r_{2m}=r_m^4,\qquad
 r_{2m+1}=X+\frac{1}{(r_m+r_{m+1})^2}.             \tag{AG.38.6}
\]
All equalities are rational-function identities, so values where
the displayed addition formula has a denominator zero are
handled by the reduced rational function. Its finite poles are
exactly the \(x\)-coordinates of nonzero points killed by \(n\).

Independently, define polynomials in \(\mathbb F_2[X]\) by
\[
 \begin{gathered}
 \psi_1=\psi_2=\psi_4=1,\qquad \psi_3=X^4+X,\\
 \psi_{2m}=\psi_m^4,\qquad
 \psi_{2m+1}=\psi_{m+2}\psi_m^3+
                       \psi_{m-1}\psi_{m+1}^3\quad(m\ge2).
 \end{gathered}                                                 \tag{AG.38.7}
\]
The program \path{primitive_parameter_rational_check.py}
constructs (AG.38.6) using reduced numerator--denominator pairs,
without using (AG.38.7) for that construction. For both required
integers it verifies that the reduced denominator is exactly
\(\psi_n^2\). Thus no general division-polynomial formula in
characteristic different from two is being imported. The checks are
\[
\begin{array}{c|r|r|r}
 n&\deg\psi_n&\deg\operatorname{num}(r_n)&
                         \deg\operatorname{den}(r_n)\\\hline
 191&18240&36481&36480\\
 229&26220&52441&52440
\end{array}
\]
and the reduced numerators are coprime to their denominators.

Finally, exact polynomial arithmetic over the explicitly defined
field \(k\) gives polynomials \(U_n,V_n\in k[u]\) satisfying
\[
       U_n(u)\psi_n(A(u))+V_n(u)\psi_n(B(u))=1
                \qquad(n=191,229).                            \tag{AG.38.8}
\]
The producer \path{primitive_parameter_gcd.py} evaluates
(AG.38.7) directly at the quadratic polynomials (AG.38.4), computes
the extended gcd, and checks the whole identity (AG.38.8) by
polynomial multiplication. The degrees of the two evaluated
polynomials are \(36480\) and \(52440\), respectively; in
each case both Bezout coefficients have degree four less.
All four polynomials for each integer are retained as binary
coefficient arrays. Their encoding, degrees, and SHA-256 hashes
are specified in \path{primitive_parameter_gcd.json}:
coefficients are in ascending degree, each represented by an
unsigned 64-bit little-endian integer whose bit \(i\) is the
coefficient of \(t^i\). The recorded runtime dependency is
\texttt{python-flint==0.9.0}. This is a complete exact identity
certificate, not a numerical approximation or a random-point test.
An independently written loader, \texttt{A6\_individual\_curve\_verify.py},
also verifies every coefficient-array hash, degree, and encoding,
reconstructs both compositions independently, and checks both full
Bezout identities; its retained run exits successfully.

If a finite \(Q\in C(\bar k)\) had \([Q-P]\ne0\) killed
by \(n=191\) or \(229\), its image under \(\alpha\) would
be killed by \(n\) in both elliptic factors. By (AG.38.5)--
(AG.38.7), the finite value \(u(Q)\) would be a common zero of
\(\psi_n(A(u))\) and \(\psi_n(B(u))\), contradicting
(AG.38.8). The remaining point \(P\) represents the origin
and cannot supply a primitive character. Therefore neither
remaining conductor occurs, over any finite extension or over
\(\bar k\), and the claimed assertion follows.

Finally \(J_C[2^\infty](\bar k)=0\). A collision would already
give a two-primary class \([Q-P']\) on \(C\), for the two
colliding points \(P',Q\), by R1. This class is zero, and the
Abel map of a positive-genus curve is injective on geometric
points. Hence \(Q=P'\), proving the stated A1 assertion.
\end{proof}
